\section{Convergence Analysis}
\label{sec:convergence}

This section establishes the convergence properties of the MPECSS algorithm. We prove that the sequence of iterates generated by Phases I and II converges to feasible points satisfying progressively stronger stationarity conditions, with Phase II guided by the adaptive path-following strategy to target B-stationary points. The analysis proceeds in three parts: first, we state the underlying assumptions; second, we establish convergence of the base regularization sequence; third, we characterize how the adaptive $t$-update and sign test guide the algorithm toward B-stationary solutions.

\subsection{Assumptions and Convergence Framework}
\label{sec:assumptions}

Throughout this section, we make the following standard assumptions, which are common in the MPEC literature~\cite{LuoPangRalph1996,ScheelScholtes2000,Scholtes2001}:

\begin{assumption}[Smoothness and Boundedness]
\label{ass:smoothness}
The functions $f$, $g_i$ ($i=1,\ldots,m$), $h_j$ ($j=1,\ldots,p$), $G_\ell$, and $H_\ell$ ($\ell=1,\ldots,q$) are twice continuously differentiable. The feasible set
\[
  \mathcal{F} = \{x \in \mathbb{R}^n : g_i(x) \leq 0, h_j(x) = 0, G_\ell(x) \geq 0, H_\ell(x) \geq 0, G_\ell(x) H_\ell(x) = 0\}
\]
is non-empty and bounded, or at least contains the level set $\{x : f(x) \leq f(x_0)\}$ under control.
\end{assumption}

\begin{assumption}[Constraint Qualification]
\label{ass:cq}
At the (unknown) solution $x^*$, MPEC-LICQ (Definition~\ref{def:mpec_licq}) or MPEC-MFCQ (Definition~\ref{def:mpec_mfcq}) holds. This ensures that the Lagrange multipliers in the stationarity system are well-defined (or unique under MPEC-LICQ).
\end{assumption}

\begin{assumption}[NLP Solver Accuracy]
\label{ass:nlp_accuracy}
At each outer iteration $k$, the solution $x^k$ of the subproblem $\text{NLP}(t_k)$ satisfies: (1)~the KKT conditions with tolerance $\varepsilon_k > 0$, and (2) the feasibility constraints are satisfied exactly or to machine precision. The tolerance $\varepsilon_k$ is progressively reduced as $t_k$ decreases (e.g., $\varepsilon_k = O(t_k)$).
\end{assumption}

\begin{remark}
Assumption~\ref{ass:nlp_accuracy} reflects the practical implementation
where IPOPT~\citep{WachterBiegler2006} is employed as the default NLP
engine. The released code does \emph{not} implement a proportional
inner-tolerance schedule of the form $\varepsilon_k = O(t_k)$; instead,
the benchmark configuration uses a fixed inner NLP tolerance inside each
subproblem solve, together with the outer stopping tests and
post-processing logic described in Section~\ref{sec:phase_iii}. The
analysis below therefore relies only on the standard requirement that
the regularized NLPs are solved to sufficiently accurate approximate
KKT points.
\end{remark}

\subsection{Convergence of the Regularization Sequence}
\label{sec:conv_reg_seq}

The foundation of MPECSS is Scholtes's regularization scheme~\cite{Scholtes2001}, which has been extensively studied in the literature~\cite{HoheiselKanzowSchwartz2013,RalphWright2004,SteffensenUlbrich2010}. The following theorem, originally due to Scholtes~\citep{Scholtes2001}, establishes that limit points of the regularized sequence are stationary for the original MPEC.

\begin{theorem}[Scholtes's Regularization Convergence~\citep{Scholtes2001}]
\label{thm:scholtes}
Let $\{x^k\}$ be a sequence of local minimizers of the regularized subproblems $\text{NLP}(t_k)$ with regularization parameters $t_k > 0$, $t_k \to 0$. If:
\begin{enumerate}
  \item[(i)] the sequence $\{x^k\}$ is bounded, or has a convergent subsequence,
  \item[(ii)] MPEC-LICQ (Definition~\ref{def:mpec_licq}) holds at the limit point $x^*$,
\end{enumerate}
then any limit point $x^*$ of $\{x^k\}$ is feasible for the original MPEC~\eqref{eq:general_mpec} and is a C-stationary point (Definition~\ref{def:mpec_stationarity}).
\end{theorem}

\begin{proof}
The key observation is that as $t_k \to 0$, the relaxed complementarity constraint $G_\ell(x) H_\ell(x) \leq t_k$ forces the complementarity condition $G_\ell(x) H_\ell(x) = 0$ in the limit. By continuity of the constraint functions and compactness of the feasible sets, the limit $x^*$ satisfies MPEC-feasibility. C-stationarity follows from the KKT conditions of the subproblems and the sign constraints on the biactive multipliers imposed by active-set analysis. See Scholtes~\citep{Scholtes2001} and Scheel and Scholtes~\citep{ScheelScholtes2000} for full details.
\end{proof}

\begin{remark}[C-stationarity is not optimal for MPECs]
Theorem~\ref{thm:scholtes} guarantees only C-stationarity, which is strictly weaker than B-stationarity in the hierarchy~\eqref{eq:stat_hierarchy}. Indeed, standard Scholtes provides no mechanism to enforce the sign positivity conditions required for S-stationarity. This is a fundamental limitation of naive geometric reduction.
\end{remark}

\subsection{Adaptive Path-Following and Stationarity Enhancement}
\label{sec:adaptive_stationarity}

The main contribution of MPECSS is to combine Scholtes's regularization with an \emph{adaptive path-following strategy} that guides the sequence toward stronger stationarity. Unlike the fixed geometric reduction $t_{k+1} = \kappa t_k$, MPECSS adjusts $t_k$ adaptively based on the complementarity residual and the multiplier sign test, aiming to keep the solution trajectory in a region where S-stationarity is more likely.

\subsubsection{Adaptive $t$-Update Strategy}
\label{sec:adaptive_t_strategy}

At each outer iteration $k$, MPECSS evaluates the implemented homotopy
complementarity residual
\begin{equation}
  r_k = \max_{\ell=1,\ldots,q} |G_\ell(x^k)H_\ell(x^k)|
  \label{eq:comp_res_def}
\end{equation}
and uses it to decide how aggressively to reduce $t_k$. The distinct
biactive residual
\[
  \eta_k = \max_{\ell=1,\ldots,q} \min(|G_\ell(x^k)|, |H_\ell(x^k)|)
\]
is used only for stationarity screening and restoration logic.
Inspired by trust-region methods~\cite{NocedalWright2006,ConnGouldToint2000},
let $\rho_k$ denote the relative improvement:
\begin{equation}
  \rho_k = \frac{r_{k-1} - r_k}{\max(r_{k-1}, \epsilon_{\text{floor}})}.
  \label{eq:relative_improvement}
\end{equation}

The implementation uses five branches, with the stagnation regime split
into two subcases:

\paragraph{Path 1}
When pathing mode is sustained (biactive set stable, $\alpha_L \le \rho_k \le \alpha_U$ with $\alpha_L = 0.2$, $\alpha_U = 5.0$, pathing counter $c \ge 2$, and $k \ge 4$) and $\rho_k > 0.5$, the algorithm applies an aggressive, \emph{superlinear} reduction:
\begin{equation}
  t_{k+1} = \kappa^2 \, t_k.
  \label{eq:superlinear_reduction}
\end{equation}

\paragraph{Path 2}
When pathing mode is sustained and $\rho_k > 0.1$, the algorithm maintains momentum with a capped reduction:
\begin{equation}
  t_{k+1} = \min(\kappa,\,0.5)\, t_k.
  \label{eq:sprint_reduction}
\end{equation}

\paragraph{Path 3}
The stagnation counter $s$ increments when $\rho_k < 0.05$ and resets otherwise. Two sub-paths handle escalating stagnation. When $s \ge 4$, an aggressive jump is performed and the counter resets:
\begin{equation}
  t_{k+1} = \kappa^2\, t_k, \qquad (s \gets 0).
  \label{eq:stagnation_jump_reduction}
\end{equation}
When $2 \le s < 4$, a milder post-stagnation reduction is applied:
\begin{equation}
  t_{k+1} = \min(\kappa,\,0.5)\, t_k.
  \label{eq:post_stagnation_reduction}
\end{equation}

\paragraph{Path 4}
In all other cases, the algorithm uses the standard geometric reduction:
\begin{equation}
  t_{k+1} = \kappa\, t_k.
  \label{eq:tiptoeing_reduction}
\end{equation}

\begin{lemma}[Adaptive $t$ Maintains Scholtes Guarantee]
\label{lem:adaptive_t_convergence}
Under Assumptions~\ref{ass:smoothness}--\ref{ass:nlp_accuracy}, if the
sequence $\{t_k\}$ generated by the implemented adaptive rule remains
positive and $t_k \to 0$, then Theorem~\ref{thm:scholtes} applies and
any limit point is C-stationary (or stronger if the subsequent
certification stage succeeds).
\end{lemma}


\subsubsection{Convergence Toward S-Stationarity via Sign Test}
\label{sec:sign_test_convergence}

The key innovation in MPECSS is to embed a \emph{multiplier sign test}
to steer the solution toward S-stationarity. At each outer iteration
$k$, after solving $\text{NLP}(t_k)$, the algorithm extracts the
converted multipliers $(\gamma^{G,k}, \gamma^{H,k})$ and checks the
sign conditions at biactive indices $\ell \in \mathcal{I}_{00}(x^k)$:
\begin{equation}
  \gamma^{G,k}_\ell \geq -\tau, \quad \gamma^{H,k}_\ell \geq -\tau,
  \label{eq:sign_test_conditions}
\end{equation}
with a tolerance chosen adaptively in the implementation as
$\tau_k=\max\{10^{-8},\sqrt{t_k}\}$. If all biactive indices satisfy
these conditions, the iterate $x^k$ is accepted as an
S-stationarity candidate. By~\cite{ScheelScholtes2000},
S-stationarity implies B-stationarity:

\begin{theorem}[S-stationarity Implies B-stationarity {\cite{ScheelScholtes2000}}]
\label{thm:s_implies_b}
Let $x^*$ be feasible for~\eqref{eq:general_mpec}. If $x^*$ is S-stationary (Definition~\ref{def:mpec_stationarity}), then $x^*$ is B-stationary (Definition~\ref{def:mpec_stationarity}). This implication holds without requiring any constraint qualification.
\end{theorem}

\begin{remark}[Practical Implementation of the Sign Test]
In practice, the sign test is a fast algebraic check performed after
each NLP solve. The biactive indices are detected adaptively with a
tolerance $\sigma_k = \max(10^{-8}, \sqrt{t_k})$, ensuring that
thresholds scale with the regularization parameter. A sign-test pass is
used to identify an S-stationarity candidate and to steer the adaptive
path; final B-stationarity is reported only after the Phase~III LICQ or
LPEC checks.
\end{remark}

\subsection{Role of Restoration and Feasibility Phase}
\label{sec:restoration_convergence}

\subsubsection{Phase I (Feasibility) and Starting Points}
\label{sec:phase_i_convergence}

Phase I solves the auxiliary problem~\eqref{eq:phase_i}:
\begin{equation}
  \text{Phase I:} \quad \min_{x} \sum_{\ell=1}^{q} (G_\ell(x) H_\ell(x))^2 \quad \text{s.t.} \quad g_i(x) \leq 0, h_j(x) = 0, G_\ell(x) \geq 0, H_\ell(x) \geq 0.
  \label{eq:phase_i_revisit}
\end{equation}
This is a well-posed smooth NLP. Convergence of Phase I guarantees that the solution $\bar{x}$ satisfies the complementarity conditions approximately ($\sum_\ell (G_\ell(\bar{x}) H_\ell(\bar{x}))^2 \approx 0$) while respecting the original constraints. This warm-start enables Phase II to begin with a good feasible point, accelerating outer convergence.

\begin{lemma}[Phase I Provides Feasible Warm-Start]
\label{lem:phase_i_warm_start}
Suppose IPOPT is applied to Phase~I~\eqref{eq:phase_i_revisit}
starting from $x_0$, and it converges to a point $\bar{x}$ satisfying
KKT conditions to reasonable accuracy. Then $\bar{x}$ is approximately
complementarity-feasible and lies in the feasible region
of~\eqref{eq:general_mpec}. When Phase~II is initialized with $\bar{x}$
and the configured regularization parameter $t_0$ (default $t_0=1$ in
the released solver, with an additional fast-forward heuristic on
near-feasible Phase~I outputs), the regularized subproblems
$\text{NLP}(t_k)$ are well-posed from the outset.
\end{lemma}

\subsubsection{Restoration and Cascading Recovery}
\label{sec:restoration_no_harm}

When the multiplier sign test fails at iteration $k$, MPECSS attempts a \emph{restoration phase} applying one of three strategies (random perturbation, directional escape, or quadratic regularizer) to move $x^k$ into a region where the sign test passes. Crucially, restoration does not violate Scholtes's convergence guarantee:

\begin{lemma}[Restoration Preserves Scholtes Convergence]
\label{lem:restoration_preservation}
Let $x^k$ be the solution of $\text{NLP}(t_k)$ with multipliers that fail the sign test. If a restoration strategy produces a point $x^k_{\text{restored}}$ in the feasible region that also solves $\text{NLP}(t_k)$ (or approximately), then the sequence $\{x^k_{\text{restored}}\}$ satisfies the same asymptotic convergence as the Scholtes sequence. If all restoration strategies fail and the algorithm continues with the standard parameter update $t_{k+1} = \kappa t_k$, the algorithm reverts to the base Scholtes scheme and still converges to at least C-stationarity.
\end{lemma}


\subsection{B-Stationarity Certification and LPEC Post-Solve}
\label{sec:bstat_certification}

After Phase II converges (or reaches a maximum iteration count),
MPECSS executes a \emph{B-stationarity certification phase} (Phase III)
to verify that the final iterate is B-stationary, especially when
MPEC-LICQ may not hold.

\subsubsection{MPEC-LICQ Check and Direct Certification}
\label{sec:licq_check}

When the sign test is passed at the final iterate $x^*$ and MPEC-LICQ
holds at $x^*$, the algorithm certifies $x^*$ as B-stationary by the
equivalence between S- and B-stationarity under LICQ. If MPEC-LICQ
does not hold or is uncertain, MPECSS solves an LPEC to provide a
fall-back geometric verification:

\begin{definition}[LPEC for B-Stationarity~\citep{LuoPangRalph1996,FletcherLeyfferRalphScholtes2006,NurkanovicLeyffer2025}]
Given a candidate $x^*$, the LPEC (Linear Program with Equilibrium Constraints) is:
\begin{equation}
\begin{aligned}
  \text{LPEC}(x^*): \quad & \min_{d \in \mathbb{R}^n} \quad && \nabla f(x^*)^\top d \\
  & \text{s.t.} \quad && \nabla g_i(x^*)^\top d \leq 0 \quad (i \in \mathcal{I}_g(x^*)), \\
  & && \nabla h_j(x^*)^\top d = 0 \quad (j = 1, \ldots, p), \\
  & && \nabla G_\ell(x^*)^\top d \geq 0, \quad \nabla H_\ell(x^*)^\top d \geq 0 \quad (\ell \in \mathcal{I}_{00}(x^*)), \\
  & && \nabla G_\ell(x^*)^\top d \text{ free}, \quad \nabla H_\ell(x^*)^\top d \text{ free} \quad (\ell \notin \mathcal{I}_{00}(x^*)), \\
  & && x^* + d \in [x_L, x_U].
\end{aligned}
\label{eq:lpec}
\end{equation}
If the optimal objective of LPEC is $\geq 0$ (i.e., no descent direction exists), then $x^*$ is B-stationary.
\end{definition}

\begin{theorem}[LPEC Certification of B-Stationarity]
\label{thm:lpec_bstat}
If $x^*$ is feasible for~\eqref{eq:general_mpec} and the LPEC~\eqref{eq:lpec} has optimal value $\geq 0$, then $x^*$ is B-stationary. Moreover, if MPEC-LICQ holds at $x^*$, B-stationarity is equivalent to S-stationarity (Definition~\ref{def:mpec_stationarity}).
\end{theorem}

\subsubsection{Stationarity Hierarchy in MPECSS Outcomes}
\label{sec:stationarity_outcomes}

The MPECSS algorithm produces one of the following outcomes for a converged problem:

\begin{enumerate}
  \item \textbf{B-Stationary (sign test pass + MPEC-LICQ):} The solution passes the multiplier sign test and MPEC-LICQ is verified. The implementation then certifies B-stationarity by S/B equivalence under LICQ.
  
  \item \textbf{B-Stationary (LPEC certified):} LPEC certification (Theorem~\ref{thm:lpec_bstat}) confirms no descent direction exists, verifying B-stationarity.
  
  \item \textbf{S-Stationary candidate (sign test pass, MPEC-LICQ uncertain):} The multiplier sign test passes, confirming an S-stationarity candidate. When MPEC-LICQ holds at $x^*$, S-stationarity is equivalent to B-stationarity~\citep{ScheelScholtes2000}; when MPEC-LICQ fails, LPEC certification (Theorem~\ref{thm:lpec_bstat}) is used to confirm or refute B-stationarity.
  
  \item \textbf{C-Stationary or weaker:} If the algorithm terminates before the sign test passes, the solution is guaranteed to be at least C-stationary by Theorem~\ref{thm:scholtes}.
\end{enumerate}

\subsection{Convergence Summary and Computational Implications}
\label{sec:convergence_summary}

\begin{theorem}[Main Convergence Result for MPECSS]
\label{thm:mpecss_convergence}
Under Assumptions~\ref{ass:smoothness}--\ref{ass:nlp_accuracy}, the MPECSS algorithm (Phases~I--III) generates a sequence $\{x^k\}$ such that:
\begin{enumerate}
  \item[(i)] The sequence is bounded (or has a limit point) in $\mathbb{R}^n$.
  \item[(ii)] Any limit point $x^*$ is feasible for the original MPEC~\eqref{eq:general_mpec}.
  \item[(iii)] If the algorithm returns a solution with status ``converged'' and sign test PASS, then $x^*$ is S-stationary; the reported B-stationarity label is then produced by the Phase~III LICQ/LPEC certification logic.
  \item[(iv)] If MPEC-LICQ (Definition~\ref{def:mpec_licq}) holds at $x^*$, then S-stationarity and B-stationarity are equivalent.
  \item[(v)] If the algorithm terminates early (e.g., maximum iterations reached) and LPEC certification passes, then $x^*$ is B-stationary.
  \item[(vi)] If the algorithm terminates without sign test pass and LPEC cannot verify B-stationarity, then $x^*$ is guaranteed to be at least C-stationary under MPEC-LICQ, by Theorem~\ref{thm:scholtes}.
\end{enumerate}
\end{theorem}


\begin{remark}[Practical Convergence Rates]
While Theorem~\ref{thm:mpecss_convergence} establishes asymptotic convergence, the rate depends on the conditioning of the MPEC and the effectiveness of the adaptive $t$-update. With the superlinear reduction (Path 1, Equation~\eqref{eq:superlinear_reduction}), the outer loop convergence can be very fast once locked on; typical runs achieve $r_k < 10^{-7}$ within 10--30 outer iterations on well-behaved problems. The restoration phase and early stagnation detection further accelerate convergence by avoiding plateaus.
\end{remark}

\subsection{Local Optimality and Constraint Qualifications}
\label{sec:local_optimality}

It is important to note that the convergence results in this section apply to \emph{feasible points} of the complementarity structure, not necessarily local minimizers of the original objective $f(x)$. Since B-stationarity is a necessary (but not sufficient) condition for local optimality~\citep{LuoPangRalph1996,ScheelScholtes2000}, MPECSS certifies that no descent direction exists within the feasible set (in the Bouligand sense). When combined with strict complementarity and higher-order analysis (not covered here), B-stationarity can be upgraded to local optimality; see~\cite{OutrataKocvaraZowe1998} for details.

\paragraph{Remark on Unbounded Feasible Sets}
If the feasible set is unbounded, Assumption~\ref{ass:smoothness} should be replaced with: the objective $f$ has a lower bound on $\mathcal{F}$, or standard growth conditions hold. Under these modifications, the convergence analysis remains valid, and any limit point of an unbounded sequence $\{x^k\}$ satisfies the same stationarity certificates.



